\documentclass{report}
\usepackage{amsmath,amssymb}
\setlength{\parindent}{0mm}
\setlength{\parskip}{1em}
\begin{document}
\begin{center}
\rule{6in}{1pt} \
{\large Paul Tsuji \\
{\bf Analysis of Shifted Laplacian Preconditioners through Chebyshev Polynomials}}

731 Oakland Ave \\ Apt 10 \\ Oakland \\ CA 94611
\\
{\tt ptsuji@gmail.com}\\
Raymond Tuminaro\end{center}


In the areas of acoustics, elasticity, and electromagnetics, the
frequency domain approach to wave propagation is very common,
particularly when the bandwidth of the signal is narrow and only a few
frequencies need to be investigated. Such computation requires the
solution of the indefinite Helmholtz problem
\begin{equation}
-\Delta u - \frac{\omega^2}{c^2} u = f
\end{equation}
where $u$ is the field of interest, $f$ is the source function, $\omega$
is the angular frequency, and $c$ is the wavespeed of the medium. For
infinite domain problems, approximations to the Sommerfeld radiation
conditions are prescribed at the boundary; for cavity or resonance
problems, Dirichlet boundary conditions are enforced.

Finite element and finite difference discretizations of the equation
above result in sparse linear systems where the matrix is indefinite and
badly conditioned; iterative methods and preconditioners which work
nicely for positive-definite elliptic problems typically do not perform
well. Traditional multigrid diverges when applied to these problems for a
few reasons. First, the indefiniteness of the operator affects the
damping
properties of smoothers like Jacobi or Gauss-Seidel; these methods are no
longer guaranteed to converge. Secondly, the $-\omega^2$ term shifts the
eigenvalues of the Laplacian, causing eigenfunctions associated with
small eigenvalues to now be oscillatory; these modes are not represented
accurately when projected onto the coarse grid.

A method which has gained popularity in the past few years is the Shifted
Laplacian preconditioner. The basic idea is to apply an approximate
inverse of the shifted problem
\begin{equation}
-\Delta u - (\alpha + i \beta)\frac{\omega^2}{c^2} u = f
\label{eq:shifted}
\end{equation}
as a preconditioner to the original problem. If $\beta$ is chosen to be
small enough, the inverse operator of the shifted problem can act as a
good preconditioner to the original problem. Adding the complex shift
term to the operator creates the affect of damping in the system, making
the problem more elliptic and more conducive to multigrid convergence. In
practice, using multigrid with the shifted problem works well as a
preconditioner for moderately sized problems.

This talk will present an analysis of Shifted Laplacian preconditioners
through the use of Chebyshev polynomial smoothers in a two-grid setting,
and will attempt to address questions regarding the choice of the
shifting parameters $\alpha$ and $\beta$. Namely, how can one choose
these parameters to strike the right balance between multigrid
convergence and accurate inversion of the original problem? Is it
possible to use different shifting parameters on different levels of the
multigrid hierarchy to improve convergence?


\end{document}
